%! TeX program = lualatex
\documentclass[../main.tex]{subfiles}
\begin{document} \section{A bit more matrix theory}

% At the end of Example~\ref{eg:age-structure-model}, we concluded that the model \(\vec{x}(t+1) = A\vec{x}(t)\) amounts to
% \[
%   \vec{x}(t) = A^{t} \vec{x}(0).
% \]
% We introduce a tool to drastically simplify this process and to pick up intuition about the model. We will be able to completely avoid matrix exponentiation.
%
% \begin{example}
%   Suppose \(A = \begin{bmatrix} 0 & 3 \\ 0.5 & 0.5 \end{bmatrix}\) models some population dynamics with \(\vec{x}(0) = \begin{bmatrix} 50 \\ 50 \end{bmatrix}\). 
%
%   This matrix \(A\) has two eigenvectors \(\vec{u} = \begin{bmatrix} 2 \\ 1 \end{bmatrix}\) and \(\vec{u} = \begin{bmatrix} -3 \\ 1 \end{bmatrix}\). We will learn how to find them later. Take for granted that
%   \[
%     A \vec{u} = \frac{3}{2} \vec{u}
%     \quad\text{and}\quad
%     A \vec{v} = - \vec{v}.
%   \]
%
%   Therefore, \(A^{t} \vec{u} = \underline{\hspace{1in}}\) and \(A^{t} \vec{v} = \underline{\hspace{1in}}\).
%
%   \blanklines{10}
%
%   % a = 40 and b = 10
%   We can express \(\vec{x}(0) = \underline{\hspace{1cm}} \begin{bmatrix} 2 \\1 \end{bmatrix} + \underline{\hspace{1cm}} \begin{bmatrix} -3 \\ 1\end{bmatrix}\) using Method~\ref{method:linear-combination} on page~\pageref{method:linear-combination}.
%
%   \blanklines{10}
%
%   This is enough theory to easily calculate \(\vec{x}(t)\) for any \(t \ge 0\).
%
%   \blanklines{10}
% \end{example}
%
% \faStar{} This week's goal is to get to finding eigenvectors and eigenvalues as quickly as we can.
% \clearpage

We need a little bit of theory to find special vectors so that the effect of \(A\vec{v}\) is sufficiently simple. \hlwarn{Edited: ``is a sufficiently simple'' \(\to\) ``is sufficiently simple''}

\begin{itemize}
  \item Add and subtract matrices by adding and subtracting entries in the same coordinates.

    For \(2\)-by-\(2\) matrices, 
    \[
      \begin{bmatrix}
        a & b \\ 
        c & d
      \end{bmatrix}
      +
      \begin{bmatrix}
        e & f \\ 
        g & h
      \end{bmatrix}
      =
      \begin{bmatrix}
        a + e & b + f \\ 
        c + g & d + h
      \end{bmatrix}.
    \]

    For \(3\)-by-\(3\) matrices,
    \[
      \begin{bmatrix}
        a_{1} & b_{1} & c_{1} \\
        d_{1} & e_{1} & f_{1} \\
        g_{1} & h_{1} & i_{1}
      \end{bmatrix}
      +
      \begin{bmatrix}
        a_{2} & b_{2} & c_{2} \\
        d_{2} & e_{2} & f_{2} \\
        g_{2} & h_{2} & i_{2}
      \end{bmatrix}
      =
      \begin{bmatrix}
        a_{1} + a_{2} & b_{1} + b_{2} & c_{1} + c_{2} \\
        d_{1} + d_{2} & e_{1} + e_{2} & f_{1} + f_{2} \\
        g_{1} + g_{2} & h_{1} + h_{2} & i_{1} + i_{2}
      \end{bmatrix}.
    \]
    \hlwarn{Edited: the second ``\(+\)'' before the result matrix \(\to\) ``\(=\)''} \hlwarn{FIXME: verify this math/content correction}

    In general, the \(i\)-th row and \(j\)-th column of \(A + B\) is \(A_{ij} + B_{ij}\).

  \item We can multiply a matrix by a constant (scalar in matrix/linear algebra lingo) by multiplying each entry in the matrix by the constant.

    \[
      \lambda 
      \begin{bmatrix}
        a & b \\ 
        c & d
      \end{bmatrix}
      = 
      \begin{bmatrix}
        \lambda a & \lambda b \\ 
        \lambda c & \lambda d
      \end{bmatrix}
      \qquad\text{and}\qquad
      \lambda 
      \begin{bmatrix}
        a & b & c \\
        d & e & f \\
        g & h & i
      \end{bmatrix}
      = 
      \begin{bmatrix}
        \lambda a & \lambda b & \lambda c \\
        \lambda d & \lambda e & \lambda f \\
        \lambda g & \lambda h & \lambda i
      \end{bmatrix}
    \]

  \item The determinant of a matrix \(A\) is a number, denoted as \(\det(A)\).

    For \(2\)-by-\(2\) matrices,
    \[
      \det 
      \begin{bmatrix}
        a & b \\
        c & d
      \end{bmatrix} 
      =
      ad - bc.
    \]

    For \(3\)-by-\(3\) matrices,
    \[
      \det 
      \begin{bmatrix}
        a & b & c \\
        e & f & g \\
        h & i & j
      \end{bmatrix} 
      =
      a \cdot \det 
      \begin{bmatrix}
        f & g \\
        i & j
      \end{bmatrix} 
      -
      b \cdot \det 
      \begin{bmatrix}
        e & g \\
        h & j
      \end{bmatrix} 
      +
      c \cdot \det 
      \begin{bmatrix}
        e & f \\
        h & i
      \end{bmatrix} 
    \]
    Notice the signs alternate.

  \item We talked about matrix multiplication.
    \[
      \begin{bmatrix}
        a_{11} & a_{12} \\
        a_{21} & a_{22}
      \end{bmatrix}
      \begin{bmatrix}
        b_{11} & b_{12} \\
        b_{21} & b_{22}
      \end{bmatrix}
      =
      \begin{bmatrix}
        a_{11}b_{11} + a_{12}b_{21} &  a_{11}b_{12} + a_{12}b_{22} \\
        a_{21}b_{11} + a_{22}b_{21} &  a_{21}b_{12} + a_{22}b_{22}
      \end{bmatrix}
    \]
    \hlwarn{Edited: fixed the product --- lower-right of the second factor ``\(v_{22}\)'' \(\to\) ``\(b_{22}\)''; and the right column of the result ``\(a_{11}b_{21}\)'' \(\to\) ``\(a_{11}b_{12}\)'', ``\(a_{21}b_{21}\)'' \(\to\) ``\(a_{21}b_{12}\)''} \hlwarn{FIXME: verify this math/content correction}
    In general, the \(i\)-th row and \(j\)-th column \hlwarn{Edited: ``\(i\)-th column'' \(\to\) ``\(j\)-th column''} \hlwarn{FIXME: verify this math/content correction} in \(AB\) is \(A_{i1}B_{1j} + A_{i2}B_{2j} + \cdots + A_{in}B_{nj}\).

  \item Lastly, we talked about linearity \hlwarn{Edited: ``linearly'' \(\to\) ``linearity''}
    \[
      A (\lambda \vec{u} + \mu \vec{v}) = \lambda A\vec{u} + \mu A \vec{v},
    \]
    where \(A\) is a matrix and \(\vec{u}, \vec{v}\) are vectors, and \(\lambda, \mu\) are constants.
\end{itemize}
\vfill{}

{\footnotesize \faExclamationTriangle{} Be aware that the overwhelming majority of matrix-related resources \hlwarn{Edited: ``resource'' \(\to\) ``resources''} on the internet is dedicated to linear algebra as a general topic. Top-ranked search results are probably far more complicated than what we need in this course. No clue whether AI gives you sufficiently simple information.}
\clearpage

Finding roots of the characteristic polynomial of a matrix is the first step in a chain of calculations to find eigenvectors.

\begin{definition}
  The characteristic polynomial for an \(n\)-by-\(n\) matrix \(A\) is \(\det(A - xI_{n})\) where \(I_{n}\) is the identity matrix.
\end{definition}

\faStar{} The quadratic formula is useful. The roots of \(ax^{2} + bx + c\) are \(\frac{-b \pm \sqrt{b^{2} - 4ac}}{2a}\).

\begin{example}
  Write down the characteristic polynomial of \(\begin{bmatrix} 2 & 0 \\ -2 & -1 \end{bmatrix}\) and find its roots.

  \blanklines{15}
\end{example}

\begin{example}
  Write down the characteristic polynomial of \(\begin{bmatrix} 1 & 0 & 2\\ 0 & 2 & 1 \\ 0 & 0 & 3 \end{bmatrix}\) and find its roots.

  \blanklines{25}
\end{example}
\clearpage
\end{document}
